\begin{abstract}
\noindent
The Erdős--Straus workbench of The Clankers (\texttt{KokunoYumeto/erdos-straus-foundation} and the Zenodo collection 10.5281/zenodo.20401937) studies the equation $4/n=1/x+1/y+1/z$. Its main documents run to more than a thousand pages:
\begin{itemize}
\item a 619-page cumulative archive (31 August 2026) and an earlier 488-page reader (7 August);
\item a 67-page and an 86-page supplement on bounded congruence transport (8 and 9 September);
\item a 124-page reader of 20 September;
\item a continuation of 17--23 September with 28 structural items.
\end{itemize}
This reader says what these documents establish, where each result is proved, and how far each was checked here.
Its core (Sections~\ref{sec:core}--\ref{sec:families}) is proved in full here: the divisor criterion for a fixed first denominator, the halved range of the least denominator, the first two residual shells and their survivors, the classical reduction modulo $840$ by five identities, the Jacobi-symbol barrier, the least-class function, and four explicit families that solve the equation for all $n$ in four arithmetic progressions, each meeting every residue class modulo $840$ that the classical reduction leaves open. The record primes of the least-class function and the finite data of the workbench's record-shell conjecture are recomputed.
Section~\ref{sec:corpus} maps the rest: every section of the 619-page archive with its page range and its bearing on the equation, and the supplements and the continuation. It also shows that one theorem of the supplement, on shear transport between adjacent shells, holds under a weaker hypothesis than stated, and that this hypothesis cannot be weakened further. Section~\ref{sec:negative} states the negative results with their exact scope.
The workbench does not claim a proof of the conjecture, and nothing here proves or disproves it.
\end{abstract}

\section*{About this reader}\phantomsection\label{sec:about}
\addcontentsline{toc}{section}{About this reader}

\paragraph{The workbench, and who made it.}
The collective author name of the workbench is \emph{The Clankers}.
The workbench credits:
\begin{itemize}
\item u/CommonCareful3149, for the mod-$107$ arithmetic-progression observation;
\item u/UmbrellaCorp\_HR, for arithmetic, positivity-code and prime-production contributions and for the central-angle and inversive-circle theorems of the archive's §6;
\item u/CivQ17, for the diagrams behind result R05;
\item u/Far-Lifeguard519, for the initiating computational investigation;
\item u/TextBackground496 and u/JGPTech, for source ideas;
\item the owner, KokunoYumeto (u/lepthymo), for the research directions and examples, for directing and developing the investigations with AI, and for the proof and source-fidelity requirements.
\end{itemize}
These credits are from \texttt{README.md}, ``Attribution and collaboration'', and \texttt{CONTRIBUTIONS.md}.
The archive records language and computational models as contributing-model provenance for particular calculations, not as authors (archive, ``Research status and provenance'').

\paragraph{What was read.}
\begin{itemize}
\item The repository at commit \texttt{d20b32e} (23 September 2026), the latest on 26 September.
\item The 488-page reader and its source, which are in the repository.
\item The 619-page archive (\texttt{00\_ERDOS\_STRAUSS\_Project\_Reader.pdf}) and its source archive (\texttt{01\_…\_Source\_Verification\_and\_Machine\_Memory.zip}), downloaded from the Zenodo record 10.5281/zenodo.22678971.
\end{itemize}
Locators have the form ``archive §$n$, p.~$m$'' for the 619-page PDF, ``supplement §$n$'' for the 86-page supplement, and file names and line numbers for the repository.
The audit notes cited as \note{21\_}, \note{43\_}, \note{43a\_}, \note{47\_} and \note{47a\_} are on the branch \texttt{claude/claude-ab-grind-20260925} of \texttt{KokunoYumeto/zeta-function-research-reader}, in \texttt{contrib/claude-ab/grind\_20260925/}.

\paragraph{How it was checked.}
\begin{itemize}
\item The core of Sections~\ref{sec:core}--\ref{sec:families} was re-derived by the author of this reader. Most of it was refereed three times as part of an earlier reader of four workbenches \cite{WorkbenchReader} (Appendix~\ref{app:verification}).
\item The new computations of this reader were made with its own scripts: the record primes of the least-class function, the deficit sets of the record shells, the four families, and the exact scope of one theorem of the supplement (Appendix~\ref{app:verification}).
\item Everything else carries a status label saying how far it was checked.
\end{itemize}
This version has not yet had an independent referee pass of its own. Appendix~\ref{app:verification} says why, and what is planned.

\paragraph{Status labels.}
\begin{itemize}
\item \textsf{Proved here}: the proof is given in full in this reader.
\item \textsf{Audited}: the workbench's proof was read in full, every step was checked, and independent checks were run.
\item \textsf{Audited (certificate)}: the finite certificate on which a proof rests was recomputed independently, and the deduction from it was read.
\item \textsf{Audited in part}: the parts named were checked; the rest is located.
\item \textsf{Reproduced}: the finite content of a statement (its examples, tables, identities or counts on a stated range) was recomputed by an independent script; its general proof was not re-derived.
\item \textsf{Located}: the statement and its proof were found and are described; nothing was checked.
\end{itemize}
The workbench itself labels statements as proved, certificate-checked, conjectural or withdrawn. Its label is quoted where it differs from the one above.
